%---------------------------------------------------------------
%  ABSTRACT & KEYWORDS
%---------------------------------------------------------------
\begin{abstract}
Classical secure file transfer protocols (SFTP) rely on RSA and
Elliptic Curve Diffie-Hellman (ECDH), both of which are vulnerable
to Shor's algorithm running on a sufficiently powerful quantum
computer. The ``Harvest Now, Decrypt Later'' (HNDL) threat demands
an immediate transition to quantum-resistant alternatives. This paper
presents \textbf{Q-SFTP}, a production-grade post-quantum secure file
transfer system that replaces all classical cryptographic primitives
with NIST-standardized algorithms: \textbf{CRYSTALS-Kyber512}
(ML-KEM, FIPS 203) for key encapsulation and
\textbf{CRYSTALS-Dilithium2} (ML-DSA, FIPS 204) for digital
signatures, combined with AES-256-GCM for symmetric encryption
and BLAKE2b-256 for file integrity verification.

Beyond the core cryptographic layer, Q-SFTP delivers a
production-ready feature set: a four-phase mutual authentication
handshake (Q-TLS), a chunked resumable transfer protocol with
cross-session state persistence, role-based access control (RBAC)
with per-user directory isolation, automatic metadata stripping
(EXIF, PDF, DOCX), malicious file protection via extension
blacklisting and magic number validation, and a tamper-evident
activity audit log. All functionality is exposed through a
modern, browser-based web GUI backed by a Flask 3.x server and
containerised with Docker. Experimental evaluation confirms
quantum-safe resistance against cryptanalysis, man-in-the-middle,
and replay attacks, while incurring approximately 0.01\% cryptographic
overhead per file transfer.
\end{abstract}

\begin{IEEEkeywords}
Post-Quantum Cryptography, CRYSTALS-Kyber, CRYSTALS-Dilithium,
Secure File Transfer, BLAKE2b, AES-256-GCM, Resumable Transfers,
Role-Based Access Control, NIST PQC, Q-TLS
\end{IEEEkeywords}
